% interface=en xxoutputxx=pdftex

% toevoegen: --background=name

\environment mstyle 

\TitlePage{506}{0}{\TeX EXEC\\explained}{}{1}

\starttext

\subject {Introduction}

\TEX\ the program is a batch text processer. You key in a document source
file in \ASCII, and \TEX\ takes this file and interprets the \TEX\ specific
typesetting instructions embedded in the text (usually preceded by \type
{\}). 

When we launch \TEX\ it first loads a so called format file, a pre||compiled
collection of macro definitions, hyphenation patterns, and sometimes font
metric data. In our case, we use the \CONTEXT\ format. 

Each time we install a new version of \CONTEXT, we have to regenerate this
format, and each time we run \TEX\ we have to specify the format to use. To
free the user from this burden of specifying and generating, we have written
\TEXEXEC, a \PERL\ script that, like \TEX, is not bound to specific hard- and
software platform. 

Just in case you wonder why the \TEXEXEC\ logo (on this cover) looks like it
is: one of the main objectives of \TEXEXEC\ is to enable users to produce
his|/|her pages as efficient as possible. So, \TEXEXEC\ is dealing with
pages. 

\subject {Running} 

Each \CONTEXT\ user interface has its own format. The next command generates
two formats, one with the English interface, also defaulting to English
typesetting, and one for the Dutch language: 

\starttyping
texexec  --make  nl en
\stoptyping

By default, the main language matches the user interface, which is normally
what we expect from a Dutch (\type {nl}), English (\type {en}) or German
(\type {de}) version. You can default to another main language and|/|or font
by changing the file \type {cont-usr.tex}, but the next alternative works
well too: 

\starttyping
texexec  --make  --language=pl,cz,sk  --bodyfont=plr  en
\stoptyping

This will generate a format file with an english user interface, while the
main language is Polish (\type {pl}) and the default body font is the Polish
alternative \type {plr} of the Computer Modern Roman (\type {cmr}). In
addition, Czech and Slovak patterns are loaded. In a similar way you can
generate a Czech version (\type {cz} and \type {csr}). 

When the appropriate formats are present, we can process a file, say \type
{test.tex}, and typeset it by calling \TEXEXEC: 

\starttyping
texexec  test
\stoptyping

\TEXEXEC\ tries to determine what interface is used. It does so by analyzing
the file. Later we will see how we can force an interface. 

When, instead of \DVI\ output, we want to produce \PDF\ code, we say: 

\starttyping
texexec --pdf  test
\stoptyping

Often one processing cycle is not enough. When no errors are encountered,
\TEXEXEC\ will call its relative \TEXUTIL. When changes are found in the
number of pages, references and/or lists, \TEXEXEC\ will process the file
again. You can prevent this by: 

\starttyping
texexec  --once  test
\stoptyping

or

\starttyping
texexec  --runs=2  test
\stoptyping

We already mentioned \PDF\ production, that can be forced by \type {--pdf}.
Specific output can be forced by: 

\starttyping
texexec  --output=dvips,acrobat  test
\stoptyping

This means that the \DVIPS\ special driver is used (which happens to be
default anyway) combined with \PDF\ instructions written in \POSTSCRIPT. 

You can force an interface on the command line by using the \type
{--interface} switch. 

\starttyping
texexec  --interface=en test
\stoptyping

This specification is seldom needed because \TEXEXEC\ has its own way of
finding out what interface to use. A rather safe way of forcing an interface
is adding a specification to the file: 

\starttyping
% interface=en tex=pdfetex output=pdftex
\stoptyping

Beware: this line should be the first line in the file! For \TEX\ this line
is simply a comment line, but \TEXEXEC\ will read it as an instruction. The
binary to be used can also be forced on the command line: 

\starttyping
texexec  --tex=pdftex  --format=plain
\stoptyping

\subject {Encodings}

In the Eastern European Countries like Poland and Czech, the normal \ASCII\
encodings are no longer sufficient. Although course, you can access the
non||latin glyphs in a font by using for instance \type {\.z}, it makes more
sense to use the slots from 127 and onwards. This not only permits a more
convenient way of keying in a document |<|in those countries people use word
processors that show the appropriate glyphs on the screen|>| but also permits
proper hyphenation. 

Instead of supporting all kind of (sometimes very rare and weird) encodings,
the default \CONTEXT\ distribution sticks to the most commonly used ones.
More details can be found in the files prefixed by \type {enco-}, \type
{lang-} and \type {font-}. 

When using the \WEBC\ distribution, you can use a mapping mechanism to map
from the document encoding to an encoding that matches the \CONTEXT\
internals, the fonts, and the hyphenation patterns. For this reason you can
inform \TEX\ on the mapping to be applied, either on the command line, or on
the first line of the file. The latter method is to be preferred, because
that way the encoding used is also properly documented in the file. Both: 

\starttyping 
%& --translate-file=cp1250pl 
\stoptyping 

\starttyping 
% --translate=cp1250pl 
\stoptyping 

are understood. The other way of informing \TEX\ is: 

\starttyping 
texexec  --translate=il2pl  somefile
\stoptyping 

In both cases, users should be aware of the fact that their files are not on
forehand portable! 

\subject {Installation} 

At start||up \TEXEXEC\ tries to locate the file \type {texexec.ini}. This
file tells \TEXEXEC\ where to find things and what binaries to use. When fed
with \type {--verbose}, \TEXEXEC\ will report the settings found there. 

In order to run \TEXEXEC\ you need to have \PERL\ installed. Installing
\PERL\ makes sense anyway. On \UNIX\ running \PERL\ scripts is natural, but
on for instance \MSWINDOWS\ you must often launch the scripts in an indirect
way. One way is simply calling \PERL: \type {perl texexec.pl}. The \FPTEX\
distribution comes with a program called \type {runperl.exe} that, when
copied to \type {texexec.exe}, takes care of the loading. Of course you also
have to copy this program to \type {texutil.exe}. In the \TETEX\ and \FPTEX\
distributions, the installation programs will take care of most issues
involved here. 

Although \TEXEXEC\ tries hard to satisfy its hunger for additional
information, sometimes you have to provide a little help. In the \CONTEXT\
distribution there is a file called \type {texexec.rme}. This file should be
copied into \type {texexec.ini}, unless of course such a file already exists.
This file is self documented, so we only list some possible settings here: 

\starttabulate[|Tl|l|Tl|]
\HL
\NC \bf variable \NC \bf meaning \NC \bf example settings \NC \NR
\HL
\NC UsedInterfaces \NC the formats generated/used \NC en,nl,de,uk \NC \NR
\NC UserInterface  \NC the default user interface \NC en          \NC \NR
\HL
\NC TeXExecutable     \NC the TeX binary to use         \NC pdfetex \NC \NR
\NC MpExecutable      \NC the MetaPost binary to use    \NC mpost   \NC \NR
\NC MpToTeXExecutable \NC the MetaPost to TeX converter \NC mpto    \NC \NR
\NC DviToMpExecutable \NC the DVI to MetaPost converter \NC dvitomp \NC \NR
\HL
\NC TeXFormatFlag \NC the format introducer         \NC \&   \NC \NR
\NC TeXVirginFlag \NC the format generation switch  \NC -ini \NC \NR
\HL
\NC TeXFormatPath  \NC fmt files                 \NC texmf/...              \NC \NR
\NC ConTeXtPath    \NC sources                   \NC texmf/tex/context/base \NC \NR
\NC SetupPath      \NC cont-sys.tex cont-usr.tex \NC texmf/tex/base/user    \NC \NR
\NC TeXScriptsPath \NC scripts                   \NC texmf/context/perltk   \NC \NR
\HL
\stoptabulate

An example of a setting is: 

\starttyping
set  TeXExecutable  to  tex
\stoptyping

You can group settings, like

\starttyping
for  tetex  set  TeXExecutable  to  pdfetex
\stoptyping

Such grouped settings are only taken into account when the shell variable is
set: 

\starttyping
set  TeXShell  to  tetex
\stoptyping

\subject {Switches}

The most important switches are mentioned already, but there are some more.
Next we will describe them all in more detail. All switches are mentioned in
full, but when issued, you can stick to a representative first part. So,
\type {--ver} is the same as \type {--verbose}. 

\switch alone

When issued, this switch makes sure that \TEXEXEC\ does as much on its own as
possible. It can for instance be used to disable \type {fmtutil} being used 
(which can be handy when tracing installation problems). 

\switch arrange 

When you specify in a file that the resulting output should be arranged,
for instance in a booklet, you must make sure that this only happens in the
last pass. This switch runs \CONTEXT\ as many times as needed to get the
references right, and then applies the last, arranging, pass. 

\switch batch 

Process the file in batchmode, that is, don't pause at errors. More on
batchmode can be found in the \TeX Book. 

\switch color

Turn on color mode. Color setting commands in the file being processed, take
preference. 

\switch convert=fileformat

Convert the input file before processing. Normally the conversion results in
a \TEX\ file. Currently supported input formats are \type {xml} and \type
{sgml}. 

\switch environment=listofnames

Load environment before processing the file. This option can be handy in
combination with input file conversion, where no environment and|/|or layout
settings are present in the file. 

\switch fast 

Run as fast as reasonable and possible without for instance spoiling time on
including graphic data. 

\switch figures=alternative 

This option can be used to generate a document with figures. The
alternatives~\type {a}, \type {b}, and~\type {c} produce different output: 

\startitemize[a,packed]
\item  a proof/correction sheet with additional information 
       about the figure
\item  a compact, paper saving sheet of figures
\item  one figure per page, with the page clipped to the 
       boundingbox 
\stopitemize

The next example shows that you can also pass an offset to be added to the
page (only important in alternative~\type {c}). 

\starttyping 
texexec  --figures=c  --paperoffset=.5cm  *.pdf *.png *.jpg 
\stoptyping 

In the process, \TEXEXEC\ uses \TEXUTIL\ to provide the list of figures. 

\switch final

Add a final run without skipping anything. This option is typically used with
fast. 

\switch format=formatfile

Normally \TEXEXEC\ prepends the \type {cont-} prefix to the format filename.
This switch can be used to specify the full name, like \type {plain}. 

\starttyping
texexec  --format=plain  --program=pdftex  somefile
\stoptyping

\switch interface=languagecode

When a file is processed, you have to use the right interface version of
\CONTEXT. This switch can be used to force \CONTEXT\ to use a certain
interface, like \type {en} (English), \type {nl} (Dutch), \type {de}
(German), \type {uk} (Brittish) etc. 

\switch language=languagecode

This switch sets the main hyphenation language, just in case this is not done
in the file itself. 

\switch listing 

A maybe unexpected switch is the one that forces \TEXEXEC\ to produce a
listing. When we have a file, say \type {readme.now}, we can typeset a
listing by invoking: 

\starttyping 
texexec --listing --pdf readme.now 
\stoptyping 

In this case, the result is available in the file \type {texexec.pdf}, but
without \type {--pdf}, you will get a \DVI\ file. 

\switch make

In order to typeset a file, \TEX\ needs a so called format file, a
pre||compiled bunch of macros. Such a format is generated when this switch is
passed. 

\switch mode=modelist

Why not typeset more than one version of one document source? The \PDFTEX\
manuals are typeset by issuing the next few commands. Watch the renaming of
the output. 

\starttyping
texexec  --pdf  --mode=A4      --result=pdftex-a  pdftex-t
texexec  --pdf  --mode=letter  --result=pdftex-l  pdftex-t
texexec  --pdf  --mode=screen  --result=pdftex-s  pdftex-t
\stoptyping

Here the \type {mode} switch tells \CONTEXT\ to obey the mode directives in
the layout specifications. 

\switch module

Documentation of \CONTEXT\ modules and related \METAPOST\ and \PERL\ modules
comes down to converting the file in a suitable format and typesetting this
source. This switch does is all. 

\switch mptex

When \TEX\ code is embedded in a \METAPOST\ file, a sequence of pre- and
postprocessing steps must take place. When set up correctly, \TEXEXEC\ does 
this annoying job for you. 

\switch noarrange 

This switch suppresses the arrangement as set up in the source file. 

\switch nomp 

Embedding \METAPOST\ code and calling for \METAPOST\ to process these files
run||time takes some time. This switch suppresses the \METAPOST\ runs. 

\switch once

Normally \TEXEXEC\ keeps on processing a file till all references are sorted
out. It therefore analyzes the files produced by \TEXUTIL\ as well as
\METAPOST\ source code generated by \CONTEXT. This option forces \TEXEXEC\ to
process a file only once. 

\switch output=driver 

One of the reasons why \TEX\ can adapt itself so easily to new developments,
is its channel to the outside world: the \type {\special} primitive. Because
each \DVI\ and \PDF\ driver has its own set of specials, we have to tell
\CONTEXT\ explicitly what specials to use. However, normally you will load
the drivers needed in the local \type {cont-sys.tex} file. 

Valid drivers directives are \type {pdftex} for native \PDFTEX\ code, \type
{dvips} (the default), \type {dvipsone} and \type {dviwindo} (the oldest
\CONTEXT\ drivers), \type {dviview} (an experimental driver) and some more. 

\switch pages=pagenumberlist

When \PDFTEX\ came around, postprocessing became \TEX's job. This switch
tells \CONTEXT\ to output only the pages as specified. Instead of a list, you
can use the keywords \type {odd} and \type {even}. 

\switch passon=string 

You can pass additional switches to the \TEX\ program used by using \type
{--passon}, like for \MIKTEX: 

\starttyping
texexec  --passon="--src"  somefile 
\stoptyping 

This tells \MIKTEX\ that you want to enrich the \DVI\ file so that editors
can relate their cursor position to a location in the \DVI\ file. Don't
forget the \type {"}, because otherwise the \type {--src} will be seen as a
switch to \TEXEXEC. 

\switch paper=key

Like the previous one, this switch is used in postprocessing. Valid keys
\type {a4a3} for A4 printed on A3, and \type {a5a4} for A5 printed on A4. The
pages are arranged as specified by the \type {print} switch. 

\switch pdf 

This is a shortcut for \type {--output=pdftex}. Because we like \PDFTEX\ and
its output and use it quite often, we introduced this dedicated switch. 

\switch pdfarrange 

Although \CONTEXT\ is pretty well able to arrange pages itself |<|think of
making A5 size booklets and those big 8~page composes pages meant for
printing books|>| it is possible to let \CONTEXT\ rearrange pages in \PDF\
files. 

\starttyping
texexec --pdfarrange --paper=a5a4 --print=up live.pdf
\stoptyping

This command makes an A5 booklet out of the famous \TEX\ live manual as
produced by \LATEX. You have to set up \TEX\ rather large (a few meg's of
memory) and the fonts should be present on the system in order to let
\PDFTEX\ include them as efficient as possible. You can influence the page
with the following keys: 

\starttabulate[|l|l|]
\NC \type{--paperoffset} \NC gap around the page     \NC \NR
\NC \type{--backspace}   \NC inner margin            \NC \NR 
\NC \type{--topspace}    \NC top/bottom margin       \NC \NR 
\NC \type{--markings}    \NC add cutmarks            \NC \NR 
\NC \type{--addempty}    \NC add empty pages after   \NC \NR 
\NC \type{--textwidth}   \NC the original text width \NC \NR
\stoptabulate

The \type {--addempty} switch takes a list of page numbers after which to
insert an empty page. This can be handy when a single sided document is
processed and you want the backside of the title page to be empty. 

When you manipulate a text that originally is typeset in single sided mode,
passing the original text width will improve the appearance of the result a
lot. 

You can pass a list of files. All files are combined into one document. This
enables the user to add for instance a title page. 

Doing more complex things is also possible, like adding additional text to
the page, when you set up a small file with layout definitions and an simple
figure insertion loop. 

You can also use this switch to combine \PDF\ files into one file: 

\starttyping 
texexec --pdfarrange --noduplex --paper=S6 file-a file-b
\stoptyping 

The result is available in the file \type {texexec.pdf}. 

\switch pdfselect 

This switch filters pages from a file. It can be used to generate a snapshot
or summary, or to isolate a page. 

\starttyping
texexec  --pdfselect  --paper=S6  --selection=1,9,14  file-1
\stoptyping 

The result is collected in the file \type {texexec.pdf}. Setting the paper
size is optional. Just like with the \type {--pdfarrange} option, you can
influence a bit the layout. 

\switch print=key

Page imposition can be done by \CONTEXT, which is normally more robust that
using a post processor, if only because \CONTEXT\ knows what (colorful)
content it is dealing with. The \type {up} key results in 2 pages per sheet
doublesided, and the \type {down} key in 2~rotated pages per sheet
doublesided. Use this switch together with the \type {paper} switch. 

\switch result=filename

When producing several typeset alternatives of one file in batch, we don't
want to overwrite the same file again and again. This switch specifies the
name of the output file. 

\switch runs=number 

When you know in advance how many runs are needed, you can overrule the
automatic number||of||runs||needed algorithm. 

\switch silent 

When you got tired of all the messages \CONTEXT\ sends you, you can silence
it by issuing this switch. 

\switch tex=programname

Once a system is set up properly, you never have to use this switch. But
occasionally, we want to run an experimental version, so here is the way to
specify the binary. 

\switch verbose 

Some of the information that \TEXEXEC\ needs, is fetched from the \type
{texexec.ini} file. This switch tells \TEXEXEC\ to write to the terminal what
it found in there. 

\switch help 

This is a rather helpful switch. Try it. You can request more help, like: 

\starttyping
texexec  --help  paper
\stoptyping

\ColofonPage

\stoptext
